\documentclass[instructions.tex]{subfiles}
\begin{document}
 
\chapter{Du Salut}

\primo Tout ce que Dieu a fait dans la création du monde et dans l'incarnation de son Verbe, il l'a fait pour sa gloire et pour le salut de l'homme. Si l'homme se conforme aux desseins de Dieu, s'il observe sa loi, Dieu lui promet un bonheur souverain dans l'éternité. Si l'homme agit contre les desseins du Créateur, il sera privé de ce bonheur, il sera malheureux pour toujours; de sorte que, sans la différente déstinée des hommes, Dieu trouvera toujours sa gloire et manifestera également ses grandeurs, sa justice, en réprouvant ceux qui lui auront été rebelles; sa miséricorde, en couronnant d'une gloire immortelle ceux qui lui auront été soumis. Or, glorifier Dieu, en observant sa loi, pour être heureux dans l'éternité et pour éviter d'y être malheureux, c'est ce qu'on appelle travailler son salut; d'où il faut conclure que le salut est l'affaire nécessaire, l'affaire importante, l'affaire propre et personnelle, l'affaire unique et essentielle de l'homme.

\secundo Quelques affaires que vous ayez, vous n'en avez donc point de plus nécessaire que votre salut, puisque vous n'êtes au monde que pour vous sauver: tout le reste n'est que vanité. Il n'est pas nécessaire que vous soyez au monde; mais il est nécessaire pour vous d'être sauvé. Dieu pouvait se dispenser de vous créer; nais il ne peut vous dispenser de travailler à votre salut. 
Le salut de l'homme est l'affaire importante puisque Dieu y a pensé de toute éternité. C'est une affaire de telle importance, que pour la faire réussir, Dieu a donné son fils, et que ce fils adorable a souffert l'agonie, les supplices et la croix. Oh! que mon salut est important, puisque, pour me sauver, un Dieu descend sur la terre, se fait victime, sue le sang et souffre la mort! 
Le salut est l'affaire propre de l'homme, tout l'avantage ou toute la perte en est pour lui. Si l'homme est sauvé, Dieu n'en est pas moins heureux; si l'homme est réprouvé, ce malheur est pour lui seul: Dieu n'est ni moins grand ni moins glorieux.

Le salut est l'unique affaire de l'homme. Si nous faisons notre salut, tout est fait, tout est gagné pour nous; si nous ne le faisons pas, tout est perdu sans ressource. 

Le salut est l'affaire essentielle de l'homme parce que ce n'est que pour cela qu'il est homme. Travailler à son salut et servir Dieu, voilà tout l'homme, dit le sage\footnote{Hoc est enim omnis homo. Eccl. 12.}.

Puisque l'homme n'est sur la terre que pour s'appliquer à son salut, c'est un plus grand désordre de voir l'homme qui ne pense pas à se sauver, que de voir les astres sans mouvemnent. L'homme qui ne travaille pas à son salut est un monstre qui doit cesser de vivre, comme le soleil devrait cesser d'être, s'il était sans activité et sans lumière.

O hommes! comprenez une bonne fois la fin pour laquelle vous êtes sur terre, et les desseins de Dieu sur vous. Si vous ne répondez pas à ses desseins adorables, vous êtes inutiles au monde: comme l'arbre infructueux, vous serez coupés et jetés au feu. Que vous reviendra-t-il de tout ce que vous faites sur la terre, si vous ne pensez pas à vous sauver ? Quel avantage tirerez-vous des créatures, si après avoir oublié votre salut, vous êtes réprouvés de votre Créateur et perdus sans ressource?

\section{Résolutions}
\primo Je graverai profondément dans mon coeur cette pensée: il faut que je me sauve, quoi qu'il m'en coûte

\secundo Et je me dirai à moi-même tous les matins en m'habillant: \og voici une journée que Dieu veut bien m'accorder encore pour le servir et travailler à mon salut: il faut donc que tout tende en moi, aujourd'hui, vers ce grand but\fg.

O mon Dieu! pénétrez-moi de l'importance et de la nécessité de mon salut, et faites-moi la grâce de le désirer, de le vouloir, et d'y travailler efficacement tous les jours de ma vie.
\end{document}
